\section{Demonstration}
\label{sec:demo}
Visitors open LabelLens at \url{https://seungbum.github.io/labellens/} (no login). A typical walk-through takes three minutes. (1)~Paste a post link or a handle, or click an example; the census bar fills as \labValidEndpoint services are asked. (2)~Compare the four lens cards: the same item may be shown as is to a logged-out visitor, with an info badge in the popular stack and with a warning in the protective stack, each with the labels that caused it. (3)~Open \emph{your stack} and toggle labelers to see which one decides the outcome. (4)~Scan the per-labeler list, where struck-through labels were retracted, faded ones expired, and dashed ones have no effect. (5)~Switch to the Network page to see how few services are reachable and how unevenly labels are spread. At the conference we run the same public build, on our laptop or the attendee's phone: attendees can inspect their own accounts, and for any account other than the examples, services without CORS are shown as \emph{does not allow browser queries}, so no collected labels about individuals are displayed. This is not a gap in the tool but the point of the census bar: attendees see on their own account the limit measured in Section~\ref{sec:eval}, where only \labBrowserPctOfEndpoint\% of declared endpoints can be read from a browser. Researchers can rebuild the full snapshot on their own machine with one command from the repository (\texttt{make measure derived local}, several hours of collection); that build is never published. The demo needs only a browser and Internet access; if the venue network blocks labeler services, the census still shows what remains reachable, and the recorded video covers the full walk-through. A two-minute video is linked from the site's About page and the repository, \url{https://github.com/seungbum/labellens}.
